\documentclass[11pt]{article}
\usepackage[margin=1in]{geometry}
\usepackage{amsmath,amssymb}
\newcommand{\finalanswer}[1]{\par\medskip\noindent\fbox{\parbox{0.94\linewidth}{\textbf{Answer:} #1}}}
\title{STAT 412 Quiz 5 --- Question 5}
\author{}\date{}
\begin{document}\maketitle
\section*{Problem}
Student heights are approximately normal with mean $172.1$ cm and standard
deviation $6.8$ cm. Suppose $200$ random samples of size $25$ are drawn, and
the sample means are recorded to the nearest tenth. Complete parts (a)--(c).

\section*{Work}
(a) For the sampling distribution of a sample mean,
\[
\mu_{\bar X}=172.1,\qquad
\sigma_{\bar X}=\frac{6.8}{\sqrt{25}}=1.36.
\]
This $1.36$ is the standard error.

(b) Since sample means are recorded to the nearest tenth and the interval is
inclusive, the recorded values $169.7$ through $173.2$ correspond to the
continuous interval $169.65$ through $173.25$. Use those boundaries:
\[
z_1=\frac{169.65-172.1}{1.36}=\frac{-2.45}{1.36}\approx-1.80,
\]
\[
z_2=\frac{173.25-172.1}{1.36}=\frac{1.15}{1.36}\approx0.85.
\]
From the z-table, $\Phi(0.85)=0.8023$ and $\Phi(-1.80)=0.0359$.
Thus
\[
P(169.7\le\bar X\le173.2)\approx0.8023-0.0359=0.7664.
\]
The expected count is
\[
200(0.7664)=153.28\approx153.
\]
(c) ``Below $169.5$'' means the recorded sample mean is less than $169.5$.
Since values are rounded to the nearest tenth, the largest recorded value below
$169.5$ is $169.4$, whose upper boundary is $169.45$. Use cutoff $169.45$:
\[
z=\frac{169.45-172.1}{1.36}=\frac{-2.65}{1.36}\approx-1.95.
\]
From the z-table, $\Phi(-1.95)=0.0256$.
Therefore
\[
P(\bar X<169.45)\approx0.0256,\qquad
200(0.0256)=5.12\approx5.
\]

\finalanswer{(a) $\mu_{\bar X}=172.1$, $\sigma_{\bar X}=1.36$;
(b) $153$ sample means; (c) $5$ sample means.}
\end{document}
